\chapter{Triển khai bằng Docker Hub}
\label{app:deployment}

Phụ lục này hướng dẫn dựng lại nền tảng từ mã nguồn công khai và các image đã
phát hành, không cần biên dịch lại từ đầu. Toàn bộ lệnh chạy trong thư mục
\texttt{implementation} của repository.

\section{Mã nguồn và yêu cầu môi trường}

Mã nguồn nằm tại \url{https://github.com/lucas-traumen/Bigdata}, thư mục
\texttt{implementation}. Bảng~\ref{tab:deploy-requirements} tóm tắt điều kiện
cần có trên máy đích.

\begin{table}[H]
\centering
\caption{Yêu cầu môi trường để chạy nền tảng.}
\label{tab:deploy-requirements}
\small
\begin{tabularx}{\textwidth}{>{\raggedright\arraybackslash}p{3.4cm} X}
\toprule
\textbf{Hạng mục} & \textbf{Yêu cầu} \\
\midrule
Docker & Docker Engine 29.x trở lên kèm Compose plugin (đã kiểm tra với Compose v5.x) \\
Bộ nhớ & Available $\geq$ 4,6~GiB cho chế độ live cùng lớp serving; chế độ batch dùng lại tối đa 2.560~MiB \\
Đĩa & $\geq$ 15~GB trống trước lần chạy đầu (image khoảng 1~GB nén, phần còn lại cho dữ liệu và checkpoint) \\
Cổng & 1883 (MQTT), 9092 (Kafka), 8000 (API), 8088 (dashboard), 5432 (PostgreSQL); đổi được qua \texttt{.env} \\
Công cụ host & \texttt{python3} cho các script drain và báo cáo tài nguyên \\
\bottomrule
\end{tabularx}
\end{table}

\section{Nhận dạng image và tính tái lập}
\label{app:image-identity}

Báo cáo không trích mã định danh (\emph{image ID}) cố định cho từng image.
Image ID là giá trị băm của cấu hình image, nên nó thay đổi mỗi lần build lại
kể cả khi mã nguồn không đổi: thời điểm tạo layer, thứ tự layer và metadata
đều tham gia vào giá trị băm. Một image ID in trong báo cáo vì vậy chỉ đúng
với đúng bản build tại thời điểm đo và sẽ không còn tồn tại sau khi image được
build lại.

Thay vào đó, image được nhận dạng bằng cặp \textbf{tag} và \textbf{digest}.
Tag là nhãn do người phát hành đặt và có thể bị trỏ lại; digest là băm nội dung
của manifest trên registry nên không đổi. Khi cần đối chiếu chính xác một phiên
chạy, digest được đọc tại chỗ:

\begin{figure}[H]
\captionof{lstlisting}{Doc digest cua mot image sau khi pull.}
\label{lst:digest}
\begin{lstlisting}[basicstyle=\ttfamily\footnotesize,frame=single]
docker image inspect trilucas/big_data:spark-1.0.0 \
  --format '{{index .RepoDigests 0}}'
\end{lstlisting}
\end{figure}

Cùng nguyên tắc đó, công cụ benchmark ghi digest của các image đang chạy vào
\texttt{run-spec.json} của từng phiên đo. Nhờ vậy bằng chứng gắn với nội dung
image thực tế, không phụ thuộc vào một ID được chép tay vào văn bản báo cáo.
Nếu một thành phần được build lại sau phiên đo, digest của nó sẽ khác và người
đọc phát hiện được ngay khi đối chiếu.

\section{Image đã phát hành}

Năm image của nền tảng được phát hành trong một repository duy nhất trên Docker
Hub, phân biệt bằng tag theo thành phần
(Bảng~\ref{tab:hub-images}). Mỗi thành phần có thêm tag
\texttt{<thành phần>-latest} trỏ tới cùng digest.

\begin{table}[H]
\centering
\caption{Image phát hành trên Docker Hub tại repository \texttt{trilucas/big\_data}.}
\label{tab:hub-images}
\small
\begin{tabularx}{\textwidth}{l l r X}
\toprule
\textbf{Thành phần} & \textbf{Tag} & \textbf{Kích thước} & \textbf{Dùng cho} \\
\midrule
spark     & \texttt{spark-1.0.0}     & 817~MB & \texttt{spark-stream}, \texttt{spark-batch} \\
backend   & \texttt{backend-1.0.0}   & 57~MB  & FastAPI serving layer \\
bridge    & \texttt{bridge-1.0.0}    & 53~MB  & Bridge MQTT sang Kafka \\
simulator & \texttt{simulator-1.0.0} & 49~MB  & Simulator đa nguồn \\
dashboard & \texttt{dashboard-1.0.0} & 21~MB  & Nginx và giao diện tĩnh \\
\midrule
\multicolumn{2}{l}{Tổng (đã nén)} & khoảng 1~GB & \\
\bottomrule
\end{tabularx}
\end{table}

Ba image nền còn lại được kéo trực tiếp từ Docker Hub và không nằm trong
repository trên: \texttt{postgres:17.11}, \texttt{apache/kafka:4.3.1} và
\texttt{eclipse-mosquitto:2.0.22}. Các image này dùng tag chính thức nên chỉ
cần pull, không cần build.

\section{Lấy mã nguồn và image}

\begin{figure}[H]
\captionof{lstlisting}{Clone repo, pull image va gan tag cuc bo.}
\label{lst:clone-pull}
\begin{lstlisting}[basicstyle=\ttfamily\footnotesize,frame=single]
git clone https://github.com/lucas-traumen/Bigdata.git
cd Bigdata/implementation

docker pull trilucas/big_data:spark-1.0.0
docker pull trilucas/big_data:backend-1.0.0
docker pull trilucas/big_data:bridge-1.0.0
docker pull trilucas/big_data:simulator-1.0.0
docker pull trilucas/big_data:dashboard-1.0.0

# compose.yaml refers to local names, so retag after pulling
docker tag trilucas/big_data:spark-1.0.0     bigdata-iot-spark:1.0.0
docker tag trilucas/big_data:backend-1.0.0   bigdata-iot-backend:1.0.0
docker tag trilucas/big_data:bridge-1.0.0    bigdata-iot-bridge:1.0.0
docker tag trilucas/big_data:simulator-1.0.0 bigdata-iot-simulator:1.0.0
docker tag trilucas/big_data:dashboard-1.0.0 bigdata-iot-dashboard:1.0.0
\end{lstlisting}
\end{figure}

Bước gắn tag là bắt buộc: \texttt{compose.yaml} tham chiếu tên cục bộ theo mẫu
\texttt{bigdata-iot-}$\langle$\texttt{thành phần}$\rangle$\texttt{:1.0.0}. Khi
image đã có sẵn, Compose dùng lại và bỏ qua bước build, nên không cần mạng tới
Maven Central và không cần thời gian biên dịch image Spark.

\section{Cấu hình}

\begin{figure}[H]
\captionof{lstlisting}{Chuan bi thu muc du lieu va file .env.}
\label{lst:prepare}
\begin{lstlisting}[basicstyle=\ttfamily\footnotesize,frame=single]
# read-only host check: docker, RAM, disk, ports, DATA_ROOT
scripts/preflight.sh

# creates DATA_ROOT subdirectories and bootstraps .env from .env.example
scripts/prepare-host.sh
\end{lstlisting}
\end{figure}

Nếu \texttt{/data} chưa ghi được, tạo thư mục bằng quyền quản trị hoặc trỏ
\texttt{DATA\_ROOT} sang đường dẫn khác trong \texttt{.env}. Các biến cần rà
trước khi chạy: \texttt{DATA\_ROOT}, năm biến cổng, \texttt{POSTGRES\_USER},
\texttt{POSTGRES\_PASSWORD} và \texttt{KAFKA\_TOPIC}. Giá trị mặc định trong
\texttt{.env.example} là cấu hình demo, không dùng cho môi trường thật.

\section{Khởi tạo dịch vụ nền}

Khi dùng image từ registry, bỏ qua \texttt{scripts/init.sh} vì script này build
lại toàn bộ image. Thay vào đó chạy lần lượt các bước khởi tạo:

\begin{figure}[H]
\captionof{lstlisting}{Khoi tao PostgreSQL, Kafka va topic.}
\label{lst:init-services}
\begin{lstlisting}[basicstyle=\ttfamily\footnotesize,frame=single]
docker compose up -d postgres
scripts/wait-for-health.sh --timeout 120 postgres

docker compose --profile live up -d kafka
scripts/wait-for-health.sh --timeout 180 kafka
docker compose --profile live run --rm --name kafka-init-once kafka-init
\end{lstlisting}
\end{figure}

PostgreSQL tự nạp schema khi khởi tạo volume lần đầu, gồm các bảng phục vụ
\texttt{sensor\_latest}, \texttt{alerts}, \texttt{processed\_events}, hai bảng
staging và schema \texttt{gold.sensor\_hourly}. Dịch vụ \texttt{kafka-init} là
container one-shot tạo topic \texttt{sensor\_raw} với ba partition.

\section{Chạy pipeline và dừng}

\begin{figure}[H]
\captionof{lstlisting}{Chay pipeline streaming va kiem tra suc khoe dich vu.}
\label{lst:run-live}
\begin{lstlisting}[basicstyle=\ttfamily\footnotesize,frame=single]
scripts/run-mode.sh live
scripts/wait-for-health.sh spark-stream backend dashboard
\end{lstlisting}
\end{figure}

Lệnh \texttt{run-mode.sh live} khởi động Mosquitto, Bridge, Kafka, Spark
streaming và năm simulator \texttt{zoneA}--\texttt{zoneE}; dashboard mở tại
\texttt{http://localhost:8088}, API tại \texttt{http://localhost:8000/health}.

\begin{figure}[H]
\captionof{lstlisting}{Batch theo gio, kiem tra va dung he thong.}
\label{lst:batch-stop}
\begin{lstlisting}[basicstyle=\ttfamily\footnotesize,frame=single]
scripts/run-mode.sh status
scripts/run-batch-hourly.sh
scripts/run-mode.sh batch --start 2026-09-15T01:00:00Z --end 2026-09-15T02:00:00Z
\end{lstlisting}
\end{figure}

Lệnh \texttt{run-batch-hourly.sh} tính cửa sổ giờ UTC vừa kết thúc, chờ khoảng
grace để event biên giới kịp vào Silver, rồi gọi job batch với file lock. Chạy
lại nhiều lần trên cùng cửa sổ là an toàn vì kết quả Gold được thay thế theo
khóa sensor--giờ--chỉ số, không cộng dồn. Để dừng toàn bộ mà vẫn giữ dữ liệu:

\begin{figure}[H]
\captionof{lstlisting}{Dung he thong ma khong xoa du lieu.}
\label{lst:stop}
\begin{lstlisting}[basicstyle=\ttfamily\footnotesize,frame=single]
docker compose --profile live stop
\end{lstlisting}
\end{figure}

Không dùng \texttt{docker compose down -v} cho thao tác thường ngày: cờ
\texttt{-v} xóa named volume của Kafka và PostgreSQL. Đường phá hủy duy nhất
được thiết kế cho mục đích đó là \texttt{scripts/reset-pipeline.sh} với cờ
tường minh \texttt{-{}-data} và \texttt{-{}-volumes}.

\section{Kiểm chứng sau khi triển khai}

\begin{figure}[H]
\captionof{lstlisting}{Kiem tra end-to-end va phuc hoi tren may dich.}
\label{lst:verify}
\begin{lstlisting}[basicstyle=\ttfamily\footnotesize,frame=single]
python3 -m unittest discover -s tests      # unit tests, no Docker required
scripts/test-e2e.sh                        # fixture 200 event qua API
scripts/test-recovery.sh                   # SIGKILL spark-stream roi restart
scripts/resource-report.sh --once          # docker stats + query progress
\end{lstlisting}
\end{figure}

Bộ kiểm thử đơn vị chạy được trên host mà không cần Docker. Hai script còn lại
cần stack đang chạy: \texttt{test-e2e.sh} kiểm tra đường đi từ simulator tới
API, quarantine và bảng alert; \texttt{test-recovery.sh} dừng đột ngột tiến
trình Spark rồi khởi động lại để xác nhận checkpoint cho phép chạy tiếp.

\section{Phương án thay thế: build từ mã nguồn}

Nếu không dùng registry, chạy \texttt{scripts/init.sh} để build toàn bộ image
tại máy đích. Phương án này cần mạng tới Docker Hub cho base image và tới
\texttt{repo1.maven.org} cho connector Kafka và driver JDBC (khoảng 13~MB, có
kiểm tra SHA256 trong \texttt{spark/install-jars.sh}), cùng thời gian build
image Spark lâu hơn đáng kể so với pull. Kết quả cuối cùng tương đương về mặt
chức năng; chỉ digest của image là khác, đúng như lý do đã nêu ở
Mục~\ref{app:image-identity}.
